\section{The factory}
\label{sec:factory}

Every part now exists. What is missing is the glue:

\begin{console}
EchoNode         needs an   Interface&
Spi              IS an      Interface,  and needs a ByteTransport&
AvrSpiTransport  IS a       ByteTransport
\end{console}

Somebody has to create the objects, in the right order, and keep them alive as long as the program
runs. That somebody should not be \code{main}, and definitely not \code{EchoNode}.

\begin{cppcode}
class Factory
{
public:
    Factory() noexcept;

    Factory(const Factory&)            = delete;
    Factory(Factory&&)                 = delete;
    Factory& operator=(const Factory&) = delete;
    Factory& operator=(Factory&&)      = delete;

    /**
     * @brief Get the CAN driver.
     *
     * @return Reference to the driver, valid for the lifetime of the factory.
     */
    [[nodiscard]] driver::can::Interface& canDriver() noexcept;

private:
    AvrSpiTransport myTransport; /**< Byte transport. Declared before the driver that uses it. */
    driver::can::Spi myDriver;   /**< The CAN driver, bound to myTransport. */
};
\end{cppcode}

\textbf{The declaration order is a requirement, not style.} Members are constructed in the order
they are declared, and \code{myDriver} takes a reference to \code{myTransport} in its constructor.
Swap the two lines and you bind a reference to an object that has not been constructed yet. The
compiler says nothing.

\textbf{The return type is \code{Interface&}, not \code{Spi&}.} That is the whole point. The
caller gets something it can use, not something it can recognise.

\textbf{No heap.} The objects are members, the factory sits on the stack in \code{main}, and
everything has a static size. On a freestanding target there is no \code{std::unique_ptr} to reach
for --- nor is one needed. Smart pointers in a factory example solve an ownership problem that does
not exist when the number of objects is known at compile time.

\section{\texorpdfstring{\code{main}}{main}}
\label{sec:main}

\begin{cppcode}
int main()
{
    Factory factory{};
    app::EchoNode node{factory.canDriver(), ResponseId};

    while (true)
    {
        node.run();
    }
}
\end{cppcode}

Five lines. \code{main} does three things and nothing more: creates the factory, creates the
application component with what the factory hands out, and runs.

The words \code{Spi}, \code{AvrSpiTransport}, \code{SPI0} and \code{PORTC} do not appear.
\textbf{If you can read your \code{main} without knowing which hardware the program runs on, you
have got the architecture right.}

\subsection{The test of the design}

Change one line in the factory:

\begin{cppcode}
private:
    driver::can::Stub myDriver;   // was: AvrSpiTransport + Spi
\end{cppcode}

Compile for the \textbf{host computer}. It should work, and the program should run ---
\code{EchoNode} echoes frames against an array in memory instead of against a CAN bus.

If it does not, the compiler error points out exactly where some line above the interface
(\file{driver/can/interface.hpp}) has learnt something about the hardware that it should not know.

It is a test you can run in ten seconds, and it says more about the quality of the architecture
than any code review.

\section{An application you can debug with}
\label{sec:debugapp}

\code{EchoNode} is good for proving the architecture and bad for debugging with: it does nothing
visible. Ahead of \kapref{ch:bringup} it is worth building one more.

\begin{narrowtable}{@{}ll@{}}
\thead{LED} & \thead{Lights when} \\ \midrule
Green & A frame has been received and acknowledged. \\
Yellow & A frame has been sent (\code{send()} returned \code{true}). \\
Red & \code{hasError()} returned \code{true}. \\
\end{narrowtable}

Plus a button that sends a frame with a known ID and a known payload.

That is a quarter of an hour's work and it pays for itself immediately at the bench: a red LED
lighting up when you connect the transceiver says something completely different from a node that
does nothing at all.

\textbf{Keep it outside \code{EchoNode}.} The LEDs belong to the application, not to the driver,
and the driver is to remain buildable for the host computer.
